\documentclass{CEAS_EuroGNC_2024}
\usepackage[utf8]{inputenc}
\usepackage[T1]{fontenc}
\usepackage{geometry}
\usepackage{graphicx}
\usepackage{float}
\usepackage{hyperref}
\usepackage[english,indonesian]{babel}

\title{Optimasi Pemrosesan Template LaTeX Menggunakan Pendekatan Berbasis AI}
\author{Academic Support Agent}
\date{\today}

\begin{document}
\maketitle

\begin{abstract}

\end{abstract}

\chapter{PENDAHULUAN}


\chapter{TINJAUAN PUSTAKA}


% Original content (commented out):
% \documentclass{CEAS_EuroGNC_2024}

\usepackage[utf8]{inputenc}

\usepackage{graphicx}
\usepackage{amsmath}
\usepackage[version=4]{mhchem}
\usepackage{siunitx}
\usepackage{array}
\usepackage{longtable}%tabulary

\setlength\LTleft{0pt} 


\title{Guidelines and Example for the Preparation of the CEAS EuroGNC Conference Papers}


% customize the following entries for the metadata/properties of the generated PDF file
\ifpdf
\hypersetup{pdfinfo={
Title={Guidelines and Example for the Preparation of the CEAS EuroGNC Conference Papers},
Author={Nicolas Fezans},
Subject={Template and Authors Instructions for CEAS EuroGNC Conferences},
Keywords={CEAS EuroGNC; Paper Preparation Guidelines; Authors Instructions; Template}
}}
\fi


\begin{document}
\maketitle

% Argument is the width that is reserved for the author names. Authors may slightly customize this value to improve the layout (e.g. to avoid unaestetic breaking of their names)
\begin{authorList}{4cm}
\addAuthor{First author}{Job Title, Business/Academic Affiliation, Department Name, ZIP-Code, City, Country. \emailAddress{first.author@affiliation.com}}
\addAuthor{Second author}{Job Title, Business/Academic Affiliation, Department Name, ZIP-Code, City, Country. \emailAddress{second.author@affiliation.com}}
\addAuthor{Third author with a very long name}{And a very long Job Title, as well as a very long Business/Academic Affiliation, and the same for the Department Name, ZIP-Code, City, and the name of his/her Country. \emailAddress{third.author@affiliation.com}}
\addAuthor{Fourth author}{Job Title, Business/Academic Affiliation, Department Name, ZIP-Code, City, Country. \emailAddress{fourth.author@affiliation.com}}
\end{authorList}


% In case you get a small identation on the first line of the abstract 
% try not starting the abstract text on a new line but directly after
% the command "\begin{abstract}" as shown below.
% (if you know how to fix this behavior, we'll be happy to fix this
% in the template)
\begin{abstract}This document serves as example for using the \LaTeX{} template for the CEAS Conference on Guidance, Navigation and Control (EuroGNC). It also contains guidelines for the authors regarding the preparation, the content, and its presentation of the papers. The paper must include an abstract, below the title and the authors (with affiliation). The name and current email address of the corresponding authors shall be provided in footnote as shown in this example. 
\end{abstract}

\keywords{CEAS EuroGNC; Paper Preparation Guidelines; Authors Instructions; Template}

\section*{Nomenclature}

{\renewcommand\arraystretch{1.0}
\noindent\begin{longtable*}{@{}l @{\quad=\quad} l@{}}
$\Psi, \Theta, \Phi$  & Euler angles (heading, pitch, roll)\\
$F$ & force\\
$a$ & acceleration\\
$d$ & damping\\
$k$ & stiffness\\
c   & chord \\
$m$ & mass\\
\end{longtable*}}

\section{Introduction}
This document is a \LaTeX{} template for preparation of papers for CEAS Conference on Guidance, Navigation and Control. If you are reading a hard copy or PDF version of this document, you can download the latest version of the template on the page \href{https://eurognc.ceas.org/author_instructions/}{\textcolor{blue}{\underline{https://eurognc.ceas.org/author\_instructions/}}}.

\newpage
\section{File Format and Paper Length}
In order to permit long term archival of the papers, the submitted PDF files must embed all used fonts (including in vector graphics drawings and plots). Additional information will be added to the final papers (e.g. conference name, location, dates, paper number, DOI number) and they will be provided for download in a PDF/A format. To this aim please make sure that all PDF protections are disabled from the PDF files before uploading them. Archival made by the authors or their company should ideally be based on the version containing this additional information in order to prevent the spread of copies which do not include the required information to properly cite them.

\textbf{The maximum size for the PDF file is 15\,MB}. Note papers are regularly transferred (download, email, etc.) and that it is therefore desirable to prevent large files. A typical reason for overly large file size is the unnecessary or improper use of bitmap images: either in place of some vector graphics formats or with unnecessary high resolution. In many cases a file size of 1-2\,MB can be achieved without being noticed by the readers. \textbf{All fonts must be embedded in the PDF, even if they are widespread and slightly increase the size of the PDF file.}

A page limit at 12 pages is set for the main content of the paper: between the start of the Introduction until the end of the Conclusion. Title, Authors, Abstract, and Nomenclature at the beginning of the paper are not counted in this page limit. The Acknowledgment section (if present) and the References are not counted either. The rationale for this definition that when hitting restrictions on paper length, shortening the nomenclature and references might be very tempting, but this practice should be prevented as it can quickly reduce the readability of the paper (through missing or unclear definition of the symbol used) and the deteriorate the discussion of alternative approaches (often leading to additional references).

\section{Copyright Restrictions}\label{sec:copyright_restrictions}

The publication of the submitted manuscripts in the proceedings of the conference is only possible if:
\begin{enumerate}
\item The manuscript has not been published previously.
\item The manuscript does not contain copyright-infringing material.
\item An appropriate copyright statement has been selected by the authors.
\end{enumerate}
Correspondingly authors will be required to provide statements (form will be made available to this purpose) regarding those requirements. Failure to comply with the copyright restrictions would result in the paper not being published. In order to publish the papers, the copyrights' owner, usually either the author(s) or their companies, will be required to provide a signed form for copyright transfer to CEAS or, alternatively, a license to publish and distribute the material. The corresponding copyright statement will be stamped automatically at the time of processing the final paper: no mention shall be included in the PDF file provided by the author(s).

\section{Paper Number}

A paper number will be assigned by the conference organizers to all accepted papers. This number will be used to refer to each paper in the program, to generate the paper DOI numbers, and to search for the papers in the proceedings. This paper number will be automatically stamped in the top-right corner of each paper: The authors should not add the paper number themselves. If needed, e.g. for easing the cross-referencing of papers submitted to the same edition of the conference, the final paper numbers and DOI numbers can be obtained on request to the organizing committee during the week preceding the final submission deadline. 



\section{Detailed Formatting Instructions}

The styles and formats for the CEAS Conference on Guidance, Navigation and Control papers are provided hereafter. If you are using \LaTeX{}, this template already defines all required styles to prepare your manuscript. Its use, with no modification, is highly recommended. A Microsoft Word template is also available if you prefer to use Microsoft Word. Regardless of which program you use to prepare your manuscript, please use the formatting instructions contained in this document.

These formatting instructions were developed with the aim of easing the reading of the papers on computer screens and tablets. To this aim a single column format is used with a fairly large font size (for scientific papers), slightly enlarged line heights, and combined with fairly small margins. Paragraphs are visually separated by indentation of the first line and by additional spacing between paragraphs. Whilst overall these measures tend to increase the number of pages needed for each paper, it is hoped that they will also convince a sufficient number of readers to not print the papers by providing them a better on-screen experience even on small tablets and large smartphones. The conference organizers are welcoming feedback from the authors and the readers, based on which some parameters might be reevaluated for future editions of the conference.

If you are using the electronic \LaTeX{} template to format your manuscript, the correct formatting should be automatic. The sizes provided hereafter should only be relevant for authors who are not using \LaTeX{} and who need to build their own template for their software (or format everything manually). Authors should not attempt to customize the template in any way.


\subsection{Page Format}

The paper must be prepared with the international standard A4 format as defined in the ISO~216 standard ($210\times 297$\,mm). All pages must be portrait orientation: using landscape orientation is not permitted at all. If necessary some specific content (e.g. table or picture) can be turned by 90 degrees (preferably in anti-clockwise direction) but the pages in the PDF should remain in portrait mode.

All margins are 2.0\,cm for the top and bottom margins and 1.75\,cm for the left and right margins. On the first page, an additional vertical space of 2.0\,cm is inserted between the top margin and the manuscript title: this space is reserved for the organizers to insert the conference and paper information (paper number, how to cite the paper, etc.). \textbf{It is crucial that the authors do not use this space as any content placed in the top 4.0\,cm of the first page will be masked by the conference and paper information that will be added afterward.}

In the electronic version of this template, all margins as well as other formatting is preset and should not be customized. Forcing a page break, e.g. through the command \verb|\newpage|, is however acceptable if it improves the visual layout of the paper and if, as a rule of thumb, at least 75\% of the page preceding the page break was already filled.

\subsection{Standard Text}
The default font for CEAS EuoGNC papers is Times or Times New Roman, 12-point size and with a line height of 14.4 points. For Word users: the line height has been implemented in form of a multiplicative factor of 1.15 in the Word template because this definition is more suited for the inherited properties in the styles that are based on the ``Standard'' style, whereas a fixed height of 14.4\,pt would lead to problems for these derived styles. 

The first line of every paragraph should be indented by 0.75\,cm, and all lines should be single-spaced. A paragraph vertical separation of 6\,pt should be used. As just defined the vertical separation between paragraphs are sufficient: no additional lines should be skipped between paragraphs.

Note that for easing the authors own proof reading or company internal reviews a flag called ``draft'' can be activated to create a double-spaced version of the paper. In order to use this option, by inserting the option \verb|[draft]| on the first line where the documentclass (here \verb|EuroGNC|) is defined with the command: \verb|\documentclass[draft]{CEAS_EuroGNC}|. The option \verb|draft| is intended for the authors own proof reading but should not be used for the preparation of the manuscripts submitted to the conference.





\subsection{Paper Title}

The title of your paper should be typed in bold, 20-point type, with capital and lower-case letters, and centered at the top of the page. 

Titles should ideally be both concise and precise. The title must reflect the content of the final paper: if over the course of the preparation of the final paper, modifications of the title appear useful to better reflect the final content of the paper, authors should contact the conference organizers with their submission or paper number, the old title, the proposed new title as well as a short justification for the requested change. Authors should avoid asking repeatedly for changes of the paper titles as this involves unnecessary work for the organizers.


\subsection{Authors and Affiliations}

The names of the authors, their job title, organization/affiliation, department name, ZIP-code, city, (whenever useful province/state) and country should follow on separate lines below the title. The author names should be in bold, 12-pt type font on the left and the remaining information on the right side in regular, 11-point type font. The \LaTeX{} template provides a dedicated command for adding authors:

\begin{quote}\texttt{\textbackslash addAuthor\{First name(s) and last name of the first author\}\{Job\newline  Title, Business/Academic Affiliation, Department Name, ZIP-Code,\newline City, (State/Province), Country.\newline \textbackslash emailAddress\{first.author@affiliation.com\}\}}\end{quote}

The width of the first column (where the author names are) should be set at 4\,cm. The author may slightly change this value if 4\,cm would cause a particularly unaesthetic layout, e.g. in the way the name of an author's name is split. The distance between the text of both columns (left the author names, right their job title and affiliation) should be 4\,mm.

\subsection{Abstract}

An abstract is mandatory and it should appear below the authors. It should typically be between 200 and 400 words long and not an introduction for the paper! The abstract should indicate in summary form the topic, the main objectives, and the key facts and conclusions of the paper. The abstract should be understandable to any reader with reasonable background in the areas covered by the scope of the conference, not just by the specialists of the techniques used in the paper. Based on the title and of the abstract, the reader should be able to know whether reading the paper is likely to be a good investment of his/her time. The abstract should be bold and expand on the whole text width.

\subsection{Keywords}

Keywords are optional. Authors can include up to five keywords (each being possibly composed of more than one word), separated by semicolons ``;''. When used, keywords must be listed right after the abstract. A new paragraph without indentation starting with term ``\textbf{Keywords:}'' in bold 11-point font introduces the semicolon-separated list of keywords in non-bold 11-point font. Semicolons within keywords should be avoided: if absolutely necessary, semicolons can be used provided that the keywords are additionally delimited by opening (``) and closing ('') quotes. Note that a \LaTeX{} command called \verb|\keywords| is available and used in this example to introduced the list of keywords.



\subsection{Nomenclature and Units}
Papers with many symbols may benefit from a nomenclature list that defines all symbols, inserted between the abstract (or the keywords if used) and the introduction. If one is used, it must contain all the symbols used in the manuscript, and the definitions do not need to be repeated in the text. In all cases, identify the symbols used if they are not widely recognized in the profession. The standard notations recognized in the profession should preferably be used (see for instance the current versions of the ISO\,1151 standard series for flight mechanical quantities \cite{ISO1151-1:1988,ISO1151-2:1985,ISO1151-2:1985_ADD1:1987,ISO1151-3:1989,ISO1151-3:1989_COR1:1996,ISO1151-4:1994,ISO1151-5:1987,ISO1151-6:1982,ISO1151-7:1985,ISO1151-8:1992,ISO1151-9:1993,ISO1151-9:1993_AMD1:1998}). Define abbreviations and acronyms in the text, not in the nomenclature (see section~\ref{sec:abbrv_accron}).

Symbols representing a physical quantity can be defined without providing a unit, e.g. $m$ can represent the mass of an object regardless of its unit and equations like $F=m a$ can be discussed and understood at a conceptual or symbolic level. As soon as numerical values are discussed or provided, the unit must be provided either in the nomenclature or wherever numerical values are provided (e.g. table, figure, text).

\textbf{It is strongly encouraged to use SI units throughout the paper}. Applications that were defined with other units (e.g. US Eng.) should provide wherever possible the corresponding SI conversions in addition to the original values. When some particular non-SI units are predominantly used, e.g. for altitude and airspeed in aviation, these units are tolerated especially when this results in simpler numerical values (e.g. 37\,000\,ft versus 11\,277.6\,m). Please also follow the instructions related to units in figures and figure labels in section~\ref{sec:fig_graph_tables}.




\subsection{Abbreviations and Acronyms}\label{sec:abbrv_accron}
Define abbreviations and acronyms the first time they are used in the text, even after they have already been defined in the abstract. Very common abbreviations in aerospace such as NASA, ESA, DLR, ONERA, CNES, CEAS, AIAA, SI, AC, UAV, RPAS, etc. do not have to be listed in the nomenclature, but they can be defined at first use. Do not use abbreviations in the title unless they are unavoidable (for instance, ``CEAS EuroGNC'' in the title of this article).


\subsection{Section Headings}

Level-1 headings (command \verb|\section| with LaTeX) are bold 17-point font, flush left, and numbered with arabic numerals (1 2 3 etc.). An additional spacing of 18\,pt should be used before and after level 1 headings (i.e. typically 24\,pt separation from normal text).

Level-2 headings (command \verb|\subsection| with LaTeX) are bold 14-point font, flush left, and numbered based on the number of the current level 1 heading and the level 2 heading number separated with a point (1.1  then 1.2 etc. for subsections of section 1). An additional spacing of 12\,pt should be used before and after level 2 headings (i.e. typically 18\,pt separation from normal text).

Level-3 headings (command \verb|\subsubsection| with LaTeX) are italic 12-point font, flush left, and numbered following the same logic than for subsections (1.1.1 then 1.1.2 then 1.1.3 etc.). An additional spacing of 12\,pt should be used before (but not after) level 3 headings (i.e. typically 18\,pt separation from normal text as between two subsequent paragraphs).

Level-4 heading (command \verb|\paragraph| with LaTeX) are normal shape 12-point font, flush left, and not numbered.An additional spacing of 12\,pt should be used before (but not after) level 4 headings (i.e. typically 18\,pt separation from normal text as between two subsequent paragraphs).

When followed by normal text, the vertical separation between level 1, 2, 3, or 4 headings and the normal text should be 6\,pt.





\subsection{Equations}


Equations are numbered consecutively, with equation numbers in parentheses flush right, as in Eq.~\eqref{eq:implicit_model_following_LQ_criterion}. Insert a blank line above and below the equation. To insert an equation into the \LaTeX{} document, use the \verb|\begin{equation}...\end{equation}| command environment. When related equations are written on successive lines and when possible, these equations should be aligned on the most suited place to visually show the relationship between them. Make sure that the symbols in your equation are defined before the equation appears (possibly in the nomenclature), or immediately following the equation.

Do not refer to equations just by their numbers between parentheses ``(11)'' but rather as ``equation~(11)''. The abbreviation ``Eq.~(10),'' (or ``Eqs.'' when referring to more than one scalar equation) can be used except at the beginning of a sentence. If you use these abbreviations, remain consistent throughout the paper, e.g. do not write sentences like ``\dots the criterion given in Eq.~(1) combined with the penality term from equation~(3)\dots''.

As \LaTeX{} and MS Word ``equations'' can be used to write other type of mathematical expressions (e.g. inequality, optimization problem, etc.), use the appropriate term and with no abbreviation in place of ``Eq.~''. For instance, the text could read: ``A controller $K$ minimizing the cost function $J$ defined in Eq.~(1) while satisfying the linear matrix inequality (12) and the equality constraints (13-15). Avoid using the abbreviations ``Eq.'' or ``Eqs.'' with additional qualifiers, e.g. write ``Riccati equation~(\ref{eq:ricc_can})'' rather than ``Riccati Eq.~(\ref{eq:ricc_can})''.


A few sample equations with the surrounding text are provided hereafter (in blue), formatted using the preceding instructions. 


\color{blue}
In order for the closed-loop system to follow the desired reference dynamic model $\dot{x} = A_d x + B u$, the control design problem is formulated based on the linear-quadratic implicit model following criteria $J$:

\begin{equation}\begin{split}
J &= \displaystyle\int_0^{+\infty}{(\dot{x}-A_{d}x)^T (\dot{x}-A_{d}x)} \text{dt}\\
  &= \displaystyle\int_0^{+\infty}{x^T\underbrace{(A-A_{d})^T (A-A_{d})}_{Q}x}
  + {\tilde{u}^T\underbrace{B^T B}_{R}\tilde{u}}
  + {2 x^T\underbrace{(A-A_{d})^T B}_{S}\tilde{u}}\text{ dt}.
\end{split}\label{eq:implicit_model_following_LQ_criterion}\end{equation}

The solution which minimizes the cost function $J$ is the control law $\tilde{u}=K_c\ x$ with $K_c=-R^{-1}(B^T X+S^T)$ and $X$ being the solution of the Riccati equation (\ref{eq:ricc_can}).

\begin{align}
A^T X + X A - (X B + S) R^{-1} (B^T X + S^T) + Q = 0\label{eq:ricc_can}\\
\Leftrightarrow (A^T-S R^{-1} B^T) X+ X (A-B R^{-1} S^T) - X B R^{-1} B^T X + Q - S R^{-1} S^T = 0\label{eq:ricc}
\end{align}

\color{black}

Note that in the case of equation (\ref{eq:implicit_model_following_LQ_criterion}), only one number is given to these two equations (using the \verb|split| environment) as they simply show a derivation of the same expression. Here the first one makes the idea underpinning the expression clear and the second one puts the linear-quadratic structure of this expression into relief. If the text does not refer specifically at one of these ways to write $J$, then it is not required to number each separately.



\subsection{Figures, Graphs, and Tables}\label{sec:fig_graph_tables}

Figures, graphs, and tables will be reproduced exactly as submitted. Be sure to position any figures, tables, graphs, or pictures as you want them to appear in the final paper / PDF file. Authors are responsible for ensuring that no copyright-infringing material is included (see section~\ref{sec:copyright_restrictions}).

The layout of the paper should be suited for on-screen reading of the paper: Do not place all figures and tables at the end of the paper but as close as possible to the portion of the text where they are commented and referred to. The \LaTeX{} template makes use of the \verb|hyperref| package which inserts internal hyperlinks between cross-references to figures and tables (as well as sections, equations, and references) and their exact location in the PDF file. This links are however one-way links and readers might need to scroll for going back to the text, at least unless their PDF viewer provides a ``previous view'' navigation feature. Locating figures/tables close to the corresponding portion of the text improves the on-screen reading experience.

\begin{wrapfigure}{r}{0.67\textwidth}
\vspace{-\baselineskip}\centering
\includegraphics[width=0.65\textwidth]{./graphics/longitudinal_controller_vEGNC2021Template.pdf}
\caption{Flight path angle ($\gamma$) longitudinal control law}
\label{fig:longitudinal_controller}
\end{wrapfigure}

In \LaTeX{}, do not insert your tables and figures in text boxes. As shown with Fig.~\ref{fig:longitudinal_controller} the environment \verb|wrapfigure| can be used to insert a figure on the side of the text. With Microsoft Word, depending on the version used, use the ``Insert Caption'' feature (see for instance \href{https://support.office.com/en-us/article/insert-a-caption-for-a-picture-bb74994c-7f8b-457c-be85-92233177a356}{here}) or with older versions of Microsoft Word ``frames'' and ``text boxes'' might be helpful for ensuring a more robust layout of the paper. In the \LaTeX{} template, use the ``caption'' command to type caption text. The figure title must always be placed below the figure and preceded by ``Fig.\,1'' (``Fig.'' + non-breaking half-space + figure number). A single tab shall be inserted between the figure number and the figure title. This formatting is automatic with the \verb|\caption| command of the \LaTeX{} template. A non-breaking space instead of a non-breaking half-space is acceptable. Subfigures can be used and should be numbered with the following pattern ``a)'', ``b)'', etc. followed, when necessary, by a subtitle.

Table titles follow the same rules than figure titles except that they must be placed above the table and start with ``Tab.'' instead of ``Fig.''. Make sure that the layout of the tables and their content eases its comprehension. For instance in Tab.~\ref{tab:modes_ol} the numbers were formatted such that all decimal points, exponents, and sign symbols are aligned which eases the reading. The number of digits shown should not be exaggeratedly large, e.g. do not shown experimental data with 10 significant digits if the measurement error is in the order of 1\%. 

Make sure to use very clear labels for figure axes, graphs, and table column/row-heads. These labels must use only terms that are well-defined in the paper and what is included/not included or which cases are being shown must be clear as well. Figures should have no background, borders, or outlines. The use of color for figures is allowed and recommended as it can help to make the content more appealing and readable. However, the colors (brightness and saturation levels) should be chosen such that they can be well distinguished when the paper is printed in black and white. Colors are not a substitute for minimum line thickness and line types (plain, dashed, dash-dotted, etc.) and markers. Be aware that readers might have poorly calibrated displays, printers, or projectors: the more robust the choice of colors and line styles, the better. When reusing the figures from the paper for the presentations, having made a robust choice of colors and line styles helps preventing the (too) common situation where the presenter can see and distinguish the curves on his/her laptop display but the audience cannot.

In figures and diagrams, bright text with dark background (e.g. in boxes and diagrams) should be avoided, due to the enormous ink consumption that they would cause if printed. Note that, in most cases, such choice of colors provide no real added value to the presentation of the content.




\newcommand{\ePowerTen}{\text{E}}
\newcommand{\iImag}{\bm{i}}
\begin{table}[!ht]	\centering
\caption{Eigenvalues of the open-loop system}\label{tab:modes_ol}
\begin{tabular}{|r@{}r|c|c|}\hline
   \multicolumn{2}{|c|}{Eigenvalue}                         &Damping ratio [-]       &Frequency (rad/s)   \\\hline
$-3.55\ \ePowerTen-02 \ \ $&$+\ \ 8.65\ \ePowerTen-01\  \iImag $&$4.10\ \ePowerTen-02$&$8.65\ \ePowerTen-01$\\      
$-3.55\ \ePowerTen-02 \ \ $&$-\ \ 8.65\ \ePowerTen-01\  \iImag $&$4.10\ \ePowerTen-02$&$8.65\ \ePowerTen-01$\\      
$-7.20\ \ePowerTen-03 \ \ $&$+\ \ 2.53\ \ePowerTen-01\  \iImag $&$2.84\ \ePowerTen-02$&$2.53\ \ePowerTen-01$\\      
$-7.20\ \ePowerTen-03 \ \ $&$-\ \ 2.53\ \ePowerTen-01\  \iImag $&$2.84\ \ePowerTen-02$&$2.53\ \ePowerTen-01$\\\hline
\end{tabular}
\end{table}


Captions must be descriptive and concise. Make the figures, graphs, and tables as self-contained as reasonably possible. Add any required additional description should be given in the caption, unless this would excessively increase the size of the caption. Note that figures and tables should be commented and interpreted in the text, not in the caption. 

Square brackets\, [\,]\, for delimiting units should be prevented, as square brackets are usually used for physical dimensions (e.g. Length, Mass, etc.). Units can be introduced in one of the following ways:\begin{itemize}
\item using a comma and/or the word ``in'', e.g. ``$m$ in kg'', ``$m$, in kg'', or ``$m$, kg''
\item using round parentheses and possibly with the word ``in'', e.g. ``$m$ (kg)'' or ``$m$ (in kg)''.
\end{itemize}
In order to avoid confusions, the way units are introduced can be adapted between different figures or tables of the paper, but it should be the same within each figure and table.

Note that units can be used and are often needed for dimensionless quantities. A typical example are angles, which have no physical dimension but can be expressed with various units, e.g. in radian, degrees, or gradians. Whilst most authors precise whether they use radians or degrees for angle numerical values, too often derived quantities are less consistently used. For instance, an angular frequency in `rad/s' should never appear with the unit `1/s' or `-/s'.

Note that unit symbols should never be written with the mathematical font or in italic/slanted shape. This way the symbol $m$ (mass) can be distinguished from the unit m (meters). Subscripts or superscripts consisting of an abbreviation should be written with the normal text font type. For instance a ``reference acceleration'' would be written $a_\text{ref}$ and not $a_{ref}$.

\textbf{All text in figures (label and any other text) must be legible and no smaller than 10-point type.} This corresponds to the size of the subscripts ``ref'' and ``sat'' in Fig.~\ref{fig:longitudinal_controller}.

Vector graphics should preferably be used over bitmap images, especially for graphs, plots, block diagrams. Block diagrams extracted from a Simulink model (or any simular graphical tool) are often hardly readable and authors should take the time to draw better block diagrams and export them as vector graphics and with sufficiently large font, line thickness, and arrow-head size instead of making a copy/paste from their Simulink model or subsystem.

\newpage
\subsection{Quotes}
Extended quotes, such as the following example, are to be used when material being cited is longer than a few sentences, or the standard quotation format is not practical. In this \LaTeX template, the appropriate command environment is \verb|\begin{quoting}...\end{quoting}|. For authors who are not using the \LaTeX{} template, extended quotes are to be in the same font, 11-point in size, with no indentation, with reduced text width (1.0\,cm additional margins both left and right), and justified. 

\begin{quoting}
Lorem ipsum dolor sit amet, consetetur sadipscing elitr, sed diam nonumy eirmod tempor invidunt ut labore et dolore magna aliquyam erat, sed diam voluptua. At vero eos et accusam et justo duo dolores et ea rebum. Stet clita kasd gubergren, no sea takimata sanctus est Lorem ipsum dolor sit amet. Lorem ipsum dolor sit amet, consetetur sadipscing elitr, sed diam nonumy eirmod tempor invidunt ut labore et dolore magna aliquyam erat, sed diam voluptua. At vero eos et accusam et justo duo dolores et ea rebum.
\end{quoting}

The ``quoting'' environment adds an vertical spacing of 4\,pt automatically before and after the quote, leading to a total vertical separation of 10\,pt.

\subsection{References and Footnotes}
\subsubsection{References}

All references should be listed in a ``References'' section (not numbered) at the end of the paper (sorted by order of citation in the text). Bracketed numbers, e.g. ``[1]'' should be used in the ``References'' section and when citing a single paper. For multiple citations, separate reference numbers with commas \cite{ISO1151-3:1989,ISO1151-5:1987}, or use a dash to show a range \cite{ISO1151-3:1989_COR1:1996,ISO1151-4:1994,ISO1151-5:1987}. At the beginning of sentences or if the reference numbers in bracket could be confusing (e.g. in equations) references should be introduced by the word ``Ref.'' respectively ``Refs.'' followed by a non-breaking space ``\textasciitilde'' in \LaTeX{} or with the shortcut ``Ctrl+Shift+Space'' in Word.

A BibTeX style file is provided and should be used. It is derived from the well-known and freely available \texttt{unsrt} BibTeX style. One of the main modifications is that it supports the ``doi'' field for all citation types. Please provide only the DOI number in that field (e.g. just \verb|10.1007/s13272-017-0240-9| for Ref.~\cite{FezSchFis2017}) and not the full hyperlink: BibTeX will automatically create the corresponding hyperlinks.

Discussions with peers should not be cited as references but rather as footnotes.

References to publications in the English language should be preferred, but when this cannot be avoided references to publications in other languages are acceptable. In such cases, the title should be formatted as follows: ``A Translated Title in English (original title: The Original Title in the Foreign Language)''. Make sure that the capitalization rules for titles in the corresponding foreign language are used for the original title.

The following references show the formatting of different types of references:\begin{itemize}
\item See the entry for Ref.~\cite{FezSchFis2017} for an example of journal article.
\item See the entry for Ref.~\cite{zhou1996robust} for an example of book.
\item See the entry for Ref.~\cite{PhD_Thesis_Koethe} for an example of PhD thesis.
\item See the entry for Ref.~\cite{SchFezNie2019} for an example of conference paper. In the example BibTeX file, the cross-reference function of BibTeX is used. Its use is not mandatory but recommended as it permits to define the proceedings of a conference or a specific book only once, which reduces the amount of duplicate information in your \verb|.bib| file and ensures the consistency of the references when citing several papers from the same conference or chapters from the same book. Note that the parent entry (here the proceedings or book) must be located after the entries that refer to it in the \verb|.bib| file.
\item See the entry for Ref.~\cite{NASA_TM4511} for an example of technical report (here of type ``NASA Technical Memorandum'').
\end{itemize}

Note that, for aesthetic reasons, the bibliography style automatically transform all but the first letter of the title into lowercase. If this behavior must be prevented for some letters (e.g. due to a name or an acronym), this behavior can be locally disabled by using additional curly braces \{ and \}. For instance, the title of \cite{FezSchFis2017} was written: 

\begin{quote}
\verb|{In-flight remote sensing and identification of gusts, |\newline
\verb|turbulence, and wake vortices using a {D}oppler {LIDAR}}|
\end{quote}

\noindent in order to prevent the letter ``D'' of Doppler and the acronym ``LIDAR'' from being converted to lowercase. Protecting the entire title instead of the just a few letters and acronym by using double curly braces and the begining and the end of the title (i.e. \{\{ and \}\}) is not recommended.

For addresses and locations in the USA and any other country where it is commonly required to distinguish between places with the same name use the format ``City, State Name or Abbreviation, USA'', e.g. ``Atlanta, GA, USA'' and neither ``Atlanta, Georgia'' nor ``Atlanta, USA''. This applies only for the countries where it is common practice to name the the city with the corresponding state or region, e.g. simply write ``Paris, France'' and not ``Paris, Ile-de-France, France'' or similarly write ``Munich, Germany'' and not ``Munich, Bavaria, Germany''.

\subsubsection{Footnotes}
Footnotes should be placed above the 2.25\,cm margin at the bottom of the page, where they were inserted and numbered with arabic numbers starting with ``1'' and being incremented throughout the paper (not start with ``1'' again at each page). The standard \verb|\footnote| command of \LaTeX{} should directly provide the right behavior.













\section{Conclusion}
A conclusion section is not required, though it is preferred. Although a conclusion may review the main points of the paper, do not replicate the abstract as the conclusion. A conclusion might elaborate on the importance of the work or suggest applications and extensions. \textit{Note that the conclusion section is the last section of the paper that should be numbered. The appendix (if present), acknowledgment, and references should be listed without numbers.}


\section*{Appendix}

An Appendix, if needed, should appear before the acknowledgments.

\section*{Acknowledgments}
An Acknowledgments section, if used, \textbf{immediately precedes} the References. Sponsorship information and funding data are included here. 

\bibliography{CEAS_EuroGNC_sample}

\end{document}


\end{document}